% Generated draft: requires mathematical and source review.
\documentclass[11pt]{article}
\usepackage[margin=25mm]{geometry}
\usepackage{fontspec}
\setmainfont{Latin Modern Roman}
\usepackage{amsmath,amssymb,mathtools,unicode-math}
\setmathfont{Latin Modern Math}
\usepackage{enumitem,needspace,xurl,hyperref}
\setlength{\parindent}{0pt}
\setlength{\parskip}{5pt}
\setlength{\emergencystretch}{3em}
\newcommand{\N}{\mathbb{N}}
\newcommand{\Z}{\mathbb{Z}}
\newcommand{\Q}{\mathbb{Q}}
\newcommand{\R}{\mathbb{R}}
\newcommand{\C}{\mathbb{C}}
\newcommand{\F}{\mathbb{F}}
\newcommand{\A}{\mathbb{A}}
\newcommand{\PP}{\mathbb{P}}
\newcommand{\E}{\mathbb{E}}
\newcommand{\eps}{\varepsilon}
\newcommand{\dd}{\,\mathrm{d}}
\begin{document}
\section*{UnsolvedMath ID: 30005359}
\textbf{Source problem code:} catalogue entry
\subsection{Definitions and mathematical statement}
Work in the first-order language of groups \(L_{\mathrm{grp}}=\{1,\cdot,{}^{-1}\}\), with equality. All sentences below have no parameters. A group \(G\) is hyperbolic if it has a finite symmetric generating set \(S\) and a constant \(\delta\geq 0\) such that every geodesic triangle in the Cayley graph \(\operatorname{Cay}(G,S)\), with each edge of length \(1\), is \(\delta\)-thin: every point on any side lies within distance \(\delta\) of the union of the other two sides. No torsion-free hypothesis is imposed, and finite and virtually cyclic hyperbolic groups are included.

A \(\forall\exists\)-sentence means a sentence of the form
\[
\forall x_1\cdots\forall x_m\,\exists y_1\cdots\exists y_n\,
\varphi(x_1,\ldots,x_m,y_1,\ldots,y_n),
\]
where \(m,n\geq 0\) and \(\varphi\) is a quantifier-free \(L_{\mathrm{grp}}\)-formula, allowing arbitrary Boolean combinations of equations between group terms. For groups \(G,H\), write \(G\equiv_{\forall\exists}H\) when they satisfy exactly the same such sentences. Write \(G\equiv H\) when they satisfy exactly the same first-order \(L_{\mathrm{grp}}\)-sentences; this is elementary equivalence.

\textbf{Conjecture 4, as identified in the supplied report excerpt.} For every pair of hyperbolic groups \(G\) and \(H\),
\[
G\equiv_{\forall\exists}H
\quad\Longleftrightarrow\quad
G\equiv H.
\]
Thus the assertion includes groups with torsion. Elementary equivalence implies \(\forall\exists\)-equivalence directly from the definitions; the conjectural implication is the converse.

\textit{Source qualifications.} The definitions above use the usual parameter-free group-language interpretation of the notation in the supplied clean statement. The original excerpt has corrupted logical symbols, so its precise typography cannot be recovered from the supplied transcription alone. The excerpt presents the torsion-free case as established and expects a similar uniform elimination of quantifiers in the presence of torsion; it does not supply a precise statement of that elimination. It also describes a preceding, unspecified classification as a classification up to elementary equivalence conditional on this conjecture. These contextual claims are retained as supplied source assertions, not independently verified results or additional conjectures being formalized here.
\subsection{Short English statement}
The conjecture asks whether any two hyperbolic groups that agree on every parameter-free universal--existential sentence in the language of groups must agree on every first-order sentence. Torsion is allowed. The supplied record labels the problem open and describes the torsion-free case as established; these are supplied source assessments requiring independent review.
\subsection{Sources}
\begin{itemize}
\item Local UnsolvedMath registry: \texttt{data/unsolvedmath/problems.json}, with \texttt{catalog/problems.json} as the supplementary registry; supplied permanent ID \(30005359\), source problem code \texttt{catalogue entry}, title \emph{Elementary Equivalence of Hyperbolic Groups with Torsion}. This draft uses the supplied complete record; the registry files have not been independently inspected.
\item Source report: \emph{Model Theory: Combinatorics, Groups, Valued Fields and Neostability} (2023), \emph{Oberwolfach Reports}, pp.~91--149, DOI \texttt{10.4171/owr/2023/2}. Supplied source URL and literature-source link labelled ``original Oberwolfach report'': \url{https://doi.org/10.4171/owr/2023/2}. The supplied excerpt identifies the target as Conjecture~4. No individual proposer is identified in the record.
\item Supplied literature source labelled ``Sela, torsion-free hyperbolic groups'': \url{https://doi.org/10.1112/plms/pdn052}.
\item Supplied literature source labelled ``Andr\'e, elementary equivalence and hyperbolicity'': \url{https://doi.org/10.2140/gt.2020.24.1075}.
\item Literature assessment dated \textbf{2026-08-22}: this is supplied dated curation, not an independently verified assessment. It states that Sela classifies torsion-free hyperbolic groups, that the conjecture for all hyperbolic groups including torsion remains unsettled, and that later work preserves hyperbolicity under elementary equivalence without providing the claimed full classification. The supplied background expressly calls for independent verification of time-sensitive status and cited claims.
\item Additional supplied report context attributes preservation of hyperbolicity under elementary equivalence \emph{among finitely generated groups} and classification of finitely generated virtually free groups up to the equivalence denoted by corrupted symbols to earlier works labelled [1] and [2]. It also cites [9] in its torsion-free discussion. Bibliographic identifications of these numbered references are not supplied.
\end{itemize}
\end{document}
